\section{The algebraic obstruction}\label{sec:algebra}

The preceding section produced many pairs whose edge ratios have large
height, although their multiplicative defects around rectangles have
negligible height. We now show that this is impossible. The argument uses
a fixed complete bipartite pattern. Its positions will be sufficiently
independent to separate several square roots, whereas differentiation of
the rectangle relations will force a dependence between them.

Here and below a unit circle means a circle of radius exactly one in the
original real plane. For a point \(p=(x_p,y_p)\), retain the coordinates
\(z_p=x_p+iy_p\) and \(w_p=x_p-iy_p\).

\begin{proposition}[Obstruction to multiplicative rectangles]
\label{alg:obstruction}
There is no sequence with the following properties. At index \(n\), let
\(t_n\to\infty\), let \(J_n\to\infty\), and let
\(P_n\subset\R^2\) be a finite set of algebraic points with
\(\lvert P_n\rvert\le t_n^{1+o(1)}\). Let
\(G_n\subset P_n\times P_n\) be a bipartite graph, using two copies of
\(P_n\), with at least \(t_n^{2-o(1)}\) edges. To each edge \((p,r)\)
assign a nonzero algebraic number \(\lambda_{pr}\). Suppose that, for an
absolute \(c>0\),
\begin{align}
 (z_p-z_r)(w_p-w_r)
   &=2-\lambda_{pr}-\lambda_{pr}^{-1},
       \label{alg:distance}\\
 h(\lambda_{pr})&\ge cJ_n,
       \label{alg:large-height}\\
 h\left(\frac{\lambda_{pr}\lambda_{p'r'}}
                 {\lambda_{pr'}\lambda_{p'r}}\right)&=o(J_n)
       \label{alg:rectangle}
\end{align}
uniformly over edges and complete rectangles, respectively. Finally,
suppose that no edge has both endpoints on a unit circle containing at
least \(t_n^{9/10}\) points of \(P_n\). Here \(h\) is the absolute height
from \Cref{ht:height}.
\end{proposition}

We prove the proposition through a sequence of algebraic constructions.
Assume its data exist, and suppress the index when this causes no
ambiguity. At each index choose a finite normal number field \(F_n\)
containing all coordinates, ratios and \(i\). By \Cref{ht:height}, this
choice does not affect the height. In the application from
\Cref{ext:pairs}, take the field \(\mathcal N\) already used there:
every ratio is a quotient of coordinate differences.

\subsection{Negligible height as a field of constants}

Fix a nonprincipal ultrafilter \(\mathfrak u\) on \(\N\): a collection of
subsets closed under finite intersection and supersets, containing the
cofinite sets and exactly one part of every finite partition. Such a
collection is obtained by extending the cofinite filter to a maximal
proper filter. Call a set of indices \emph{large} if it belongs to
\(\mathfrak u\). Two sequences are identified when they agree on a large
set. Write
\[
 F=\prod_{n\to\mathfrak u} F_n,
 \qquad
 \mathcal U=\prod_{n\to\mathfrak u}\C,
\]
where \(F_n\) is the number field just chosen. These are characteristic-zero
fields, with \(F\subset\mathcal U\): the inverse of a nonzero class is
defined coordinatewise on its large set of nonzero entries. The field
\(\mathcal U\) is algebraically closed. Indeed, a polynomial of fixed
positive degree has that degree on a large set, and choosing one complex
root at each such index supplies a root in \(\mathcal U\). Complex
conjugation also acts coordinatewise on \(\mathcal U\).

We use limits only along \(\mathfrak u\). Thus \(a_n=o(J_n)\) means that
\(\{n:\lvert a_n\rvert<\varepsilon J_n\}\) is large for every
\(\varepsilon>0\). A graph count is called \emph{full}
if it is at least \(t_n^{2-o(1)}\) in this sense. Define
\[
 k_0=\bigl\{[u_n]\in F:h(u_n)=o(J_n)\bigr\}.
\]

\begin{lemma}\label{alg:small-height-field}
The set \(k_0\) is a subfield of \(F\), and every element of \(F\)
algebraic over \(k_0\) belongs to \(k_0\).
\end{lemma}

\begin{proof}
The height inequalities
\[
 h(uv)\le h(u)+h(v),\qquad
 h(u+v)\le h(u)+h(v)+\log 2,\qquad
 h(u^{-1})=h(u)\quad(u\ne0)
\]
prove the field assertion, since \(J_n\to\infty\). These statements do
not depend on representatives, and \(h(0)=0\).

Suppose that \(u\in F\) satisfies
\(u^m+a_{m-1}u^{m-1}+\cdots+a_0=0\) with \(a_r\in k_0\) and fixed
\(m\ge1\). At each finite index where the equation holds, the
nonarchimedean triangle inequality gives
\(\lvert u\rvert_v\le\max(1,\max_r\lvert a_r\rvert_v)\). At infinity
the same estimate holds with the additional factor \(m\): if
\(\lvert u\rvert_v>1\), divide the equation by
\(\lvert u\rvert_v^{m-1}\). Since the total weight of the infinite
places is one,
\[
 h(u)\le\sum_{r=0}^{m-1}h(a_r)+\log m.
\]
This estimate holds on a large set of indices. Its right side is
\(o(J_n)\), proving the claim.
\end{proof}

In particular, any edge ratio selected from the graphs is transcendental
over \(k_0\), by \eqref{alg:large-height}. The rectangle defects in
\eqref{alg:rectangle}, in contrast, belong to \(k_0^\times\). We next
select finitely many edges whose positions retain enough algebraic
independence to exploit this distinction.

\subsection{A complete grid with generic positions}

We first restrict the graph to plane or bounded-degree curve loci of
minimum total dimension among those retaining a full edge count. This
restriction will give uniform bounds for the vertices satisfying additional
polynomial equations, allowing us to choose a complete grid whose limiting
positions have no relations, beyond the locus equations, over the field
generated by their coefficients.

We use two standard algebraic facts. First, for every field \(K\),
\(K[Z,W]\) is a unique factorization domain
\cite[Lemma~10.120.10]{stacksUFD}. A polynomial in \(W\) with
coefficients in \(K[Z]\) is primitive if its coefficients have no
common nonconstant factor. Gauss's lemma then identifies
factorization of primitive polynomials in \(K[Z][W]\) with
factorization over \(K(Z)\): clear denominators and cancel common
prime factors of the coefficients. In particular, an irreducible
polynomial in \(K[Z,W]\) of positive \(W\)-degree remains irreducible in
\(K(Z)[W]\). Two primitive polynomials that are scalar multiples
over \(K(Z)\) differ by an element of \(K^\times\).

Second, two distinct irreducible plane curves of degrees at most
\(D\), over an algebraically closed field, have at most \(D^2\)
common points. Indeed their projective closures have the same degrees
and no common component, so B\'ezout's theorem bounds the sum of their
intersection multiplicities by \(D^2\)
\cite[Section~5.3]{fultonCurves}. Every affine intersection contributes
at least one. We will apply this both over \(\C\) at finite indices
and over the algebraically closed field \(\mathcal U\).

By a locus at an index we mean either the whole \((z,w)\)-plane or the
zero set of a nonconstant irreducible polynomial in \(\C[z,w]\) whose
degree has a fixed bound along the sequence. Its dimension is two or
one, respectively. Among pairs of such locus sequences retaining a full
number of graph edges, choose one of minimal total dimension. A minimum
exists since the pair of planes is available and the possible total
dimensions are \(2,3,4\). Finite partitions let us fix all degrees and
normalize one nonzero coefficient of each curve equation on a large
set. Denote the limiting loci in \(\mathcal U^2\) by \(V_1,V_2\), their
dimensions by \(d_1,d_2\), and the field generated over \(\Q\) by their
normalized equation coefficients and \(i\) by \(B_0\). There are finitely
many such coefficients, so \(B_0\) is countable. It is a subfield of
\(\mathcal U\), and need not be a subfield of \(F\).

Each limiting curve equation is irreducible over \(\mathcal U\).
Otherwise a factorization into positive-degree factors would give, from
finitely many coefficient sequences, a factorization at a large set of
indices. The degrees of its factors are bounded by the degree of the
equation, so there is no unbounded factorization to consider.

Each finite curve contains \(t_n^{1-o(1)}\) retained real points: divide
the full edge count by the upper bound for the other vertex set. Its
image under \((z,w)\mapsto(\overline w,\overline z)\) shares all these
points. The bounded intersection fact forces the two curves to be equal
on a large set. The limiting curve has the same symmetry. Neither
coordinate projection is constant, since a curve with constant \(z\)
would be a line \(z=c\), whose symmetric image \(w=\overline c\) is a
different line. Thus every curve equation has positive degree in both
variables.

The loci cannot both be the same unit circle. To check the exact meaning
of this assertion, suppose their limiting equations, after normalization,
were
\begin{equation}\label{alg:unit-circle}
 (z-u)(w-v)=1.
\end{equation}
Equality of its finitely many coefficients holds at a large set of
finite indices. The symmetry gives \(v=\overline u\), again an exact
coordinatewise identity on a large set. Hence the finite loci there are
the same real unit circle. Both vertex sets on it have size
\(t_n^{1-o(1)}>t_n^{9/10}\), contradicting the deletion hypothesis in
\Cref{alg:obstruction}. This reasoning concerns equality in the
ultraproduct; ordinary convergence of a radius to one would not imply
\eqref{alg:unit-circle}.

Restrict the graph to these loci and prune it to minimum degree
\(t_n^{1-o(1)}\) on both sides. Explicitly, if its initial edge and total
vertex counts are \(e_n,v_n\), repeatedly remove vertices of degree less
than \(e_n/(2v_n)\). Fewer than \(e_n/2\) edges are removed. The
remaining graph still has a full count, both vertex sets have size
\(t_n^{1+o(1)}\), and the asserted minimum degree holds.

\begin{lemma}[Uniform avoidance and a generic grid]
\label{alg:generic-grid}
For every fixed \(N,j\ge1\), one can choose a complete
\(N\)-by-\(j\) grid in the pruned graphs on a large set of indices
so that its limiting vertices
\(p_1,\ldots,p_N\in V_1\) and \(r_1,\ldots,r_j\in V_2\) have the
following property. Select both coordinates of each vertex on a plane
locus and its \(z\)-coordinate on a curve locus. All the resulting
\(Nd_1+jd_2\) coordinates are algebraically independent over \(B_0\).
The other coordinates on a curve are algebraic over these coordinates
and \(B_0\). In particular,
\begin{equation}\label{alg:joint-degree}
 \operatorname{trdeg}_{B_0(r_1,\ldots,r_j)}
       B_0(r_1,\ldots,r_j,p_1,\ldots,p_N)=Nd_1.
\end{equation}
The selected coordinates and edge ratios belong to \(F\).
\end{lemma}

\begin{proof}
First, on a side whose locus is the plane, for each fixed degree bound
\(D\) there is a \(b_D>0\) such that every nonzero polynomial of degree
at most \(D\), even with coefficients depending on \(n\), vanishes at
at most \(t_n^{1-b_D}\) pruned vertices on a large set of indices.
If this failed, choose at each index a polynomial maximizing its number
of zeros among the finite vertex set. A maximum is attained because
there are only finitely many subsets of that set. Failure of every
fixed saving means that its zero count is \(t_n^{1-o(1)}\) along
\(\mathfrak u\). One of at most \(D\) irreducible factors then contains
\(t_n^{1-o(1)}\) vertices. Their minimum degree supplies
\(t_n^{2-o(1)}\) edges after restricting this side to a bounded-degree
curve, contradicting minimality. This proves the uniform assertion.
On a curve side, a nonzero polynomial of degree \(D\) in the selected
\(z\)-coordinate has at most \(D\) zeros among the vertices, since real
points have distinct \(z=x+iy\).

Let the two pruned set sizes be \(m,n'\) and their edge density be
\(\delta\). Independent uniform vertices
\(p_1,\ldots,p_N\) have a common-neighbor fraction \(X\) on the other
side. Averaging degrees there and applying Jensen gives
\(\E X\ge\delta^N\). Therefore the number of ordered complete grids,
initially allowing repetitions, is
\begin{equation}\label{alg:grid-count}
 m^N(n')^j\E X^j
 \ge m^N(n')^j\delta^{Nj}
 =t_n^{N+j-o(1)}.
\end{equation}

Consider a fixed nonzero polynomial in the selected coordinates with
coefficients in \(B_0\). Choose representatives of its coefficients.
Its vanishing among all ordered vertex tuples, before requiring edges,
has a power saving:
\begin{equation}\label{alg:test-count}
 \#\{\text{tuples at which the test vanishes}\}
       \le t_n^{N+j-b}
\end{equation}
on a large set, for some \(b>0\) depending on the test. Here is an
induction that also handles moving coefficients. View the polynomial
as a polynomial in the last point block, and select one nonzero
coefficient polynomial in the earlier blocks. The earlier tuples
annihilating this coefficient have a power saving by induction. At
every other earlier tuple, the remaining polynomial is nonzero, with a
fixed degree bound. The uniform plane estimate, or the bounded root
count for a curve block, bounds its zeros. Multiply by the upper
bounds \(t_n^{1+o(1)}\) for the other block sizes and slightly decrease
the saving. The starting case is exactly the one-block estimates.
The finitely many coefficient identities and nonzero denominators
involved hold on a large set.

There are countably many nonzero polynomial tests over \(B_0\).
Enumerate them as \(T_1,T_2,\ldots\). For each fixed \(s\),
\eqref{alg:grid-count} and \eqref{alg:test-count}, with a finite union
bound, show that on a large set some grid avoids all the first \(s\)
tests. At index \(n\), choose the largest feasible prefix length
\(s_n\le n\), and choose a grid avoiding that prefix. Include the
validity of all representatives and denominators of a prefix in its
feasibility requirement. If no prefix is feasible, use arbitrary
number-field coordinates at that index. For every fixed \(s\), the
set where \(n\ge s\) and the first \(s\) tests are feasible is large;
hence \(s_n\ge s\) on a large set. Each fixed polynomial test is
therefore nonzero in the limiting grid. Only finite intersections of
large sets were used.

For a curve equation \(f(z,w)=0\), its positive \(w\)-degree and the
transcendence of \(z\) show that \(w\) is algebraic over \(B_0(z)\).
This proves all the transcendence assertions. The vertices and the
assigned ratios at each chosen finite grid lie in \(F_n\), so their
classes lie in \(F\).
\end{proof}

For a plane locus \(V\) over a field \(\mathcal B\), define its function
field to be \(\mathcal B(V)=\mathcal B(z,w)\). For a curve with
irreducible equation \(f\), define it to be the fraction field of
\(\mathcal B[z,w]/(f)\). A point on \(V\) is called generic over
\(\mathcal B\) if its coordinates have transcendence degree
\(\dim V\) over that field.

We record explicitly how generic coordinates identify a function field.
If \(B_0\subset\mathcal B\subset\mathcal U\) and a point on a locus
\(V\) has transcendence degree \(\dim V\) over \(\mathcal B\), then
evaluation at that point embeds \(\mathcal B(V)\) into \(\mathcal U\).
For the plane this is algebraic independence. For a curve, its equation
\(f\) is irreducible even over \(\mathcal U\), hence over
\(\mathcal B\). Its coefficients when viewed as a polynomial in \(w\)
have no common factor of positive degree in \(z\). They cannot all
vanish at an abscissa in \(\mathcal U\): a common root would give a
linear common factor over \(\mathcal U\), contradicting
irreducibility. If \(z\) were algebraic over \(\mathcal B\), the
nonzero specialized equation would make \(w\) algebraic as well.
Thus \(z\) is transcendental. Gauss's lemma makes \(f\)
irreducible over \(\mathcal B(z)\), so every polynomial relation at the
point is a multiple of \(f\). Every nonzero denominator in the function
field consequently evaluates to a nonzero element. The same argument
applies to a point after adjoining the other grid points to
\(\mathcal B\), using \Cref{alg:generic-grid}.

\subsection{Differentiation and the surviving generic point}

Take \(N=20\) and \(j=5\) in \Cref{alg:generic-grid}, and write
\(\lambda_{i\ell}\) for the edge ratios. Define
\[
 g_i=\lambda_{i1},\qquad
 h_\ell=\frac{\lambda_{11}}{\lambda_{1\ell}},\qquad
 \chi_{i\ell}=\frac{\lambda_{i\ell}\lambda_{11}}
                         {\lambda_{i1}\lambda_{1\ell}}.
\]
The rectangle condition gives \(\chi_{i\ell}\in k_0^\times\), and
\begin{equation}\label{alg:factorization}
 \lambda_{i\ell}=g_i h_\ell^{-1}\chi_{i\ell}.
\end{equation}
Repeated rows or columns would violate the algebraic independence just
proved; all vertices on either side are distinct.

\begin{lemma}\label{alg:derivation}
There is a \(k_0\)-derivation \(\partial:F\to F\) such that
\(\partial\lambda_{i\ell}\ne0\) for every grid edge.
\end{lemma}

\begin{proof}
A derivation is an additive map satisfying the product rule; being a
\(k_0\)-derivation means that it vanishes on \(k_0\). Each
\(\lambda_{i\ell}\) is transcendental over \(k_0\) by
\Cref{alg:small-height-field}. Put it in a transcendence basis of
\(F/k_0\), prescribe its derivative to be one and those of the other
basis elements to be zero, and extend by the quotient rule to the
rational function field. Since the characteristic is zero, this
derivation extends uniquely over the algebraic extension to \(F\)
\cite[Theorems~1.5 and~4.1]{conradSeparability}.
For an infinite algebraic extension, take the compatible extensions
to its finite subextensions; uniqueness makes them agree on overlaps
\cite[Section~5]{conradSeparability}.

We have thus obtained a witness derivation nonzero on each chosen
ratio. The derivations form an \(F\)-vector space. A linear combination
of the finitely many witnesses is nonzero on all ratios provided its
coefficients avoid the zeros of a product of finitely many nonzero
linear forms. This product is a nonzero polynomial over the infinite
field \(F\), so such coefficients exist.
\end{proof}

Put
\[
 A_i=\frac{\partial g_i}{g_i},\qquad
 A'_\ell=\frac{\partial h_\ell}{h_\ell},\qquad
 H_i=(\partial x_{p_i},\partial y_{p_i}),\qquad
 U_\ell=(\partial x_{r_\ell},\partial y_{r_\ell}).
\]
The derivation kills \(i\), which is algebraic over \(\Q\subset k_0\).
Define
\[
 L_\ell(z,w)=(z-z_{r_\ell})(w-w_{r_\ell}),
 \qquad R_{i\ell}=\lambda_{i\ell}^{-1}-\lambda_{i\ell}.
\]
Differentiating \eqref{alg:distance} and \eqref{alg:factorization} gives
\begin{align}
 2(p_i-r_\ell)\cdot(H_i-U_\ell)
       &=R_{i\ell}(A_i-A'_\ell),
          \label{alg:differential}\\
 R_{i\ell}^{\,2}
       &=L_\ell(p_i)\bigl(L_\ell(p_i)-4\bigr),
          \label{alg:radicand}\\
 A_i-A'_\ell&\ne0.
          \label{alg:nonzero-derivative}
\end{align}
Here the dot denotes the bilinear expression \(xu+yv\) in the ordinary
two coordinates, with no conjugation. The second identity follows by
squaring \(\lambda^{-1}-\lambda\), and the third by
\Cref{alg:derivation}.

We now form a coefficient field containing the right-hand positions and
their derivative data. Let \(K_0\) be the subfield of \(\mathcal U\)
consisting of all elements algebraic over
\[
 B_0(r_1,\ldots,r_j,U_1,\ldots,U_j,A'_1,\ldots,A'_j).
\]
This field is algebraically closed: roots in \(\mathcal U\) of
polynomials over it are still algebraic over the displayed field.
We do not adjoin \(k_0\) to \(K_0\).

\begin{lemma}\label{alg:surviving-point}
Some \(p=p_i\) has transcendence degree \(d_1\) over \(K_0\), and
evaluation identifies the function field \(K_0(V_1)\) with \(K_0(p)\).
\end{lemma}

\begin{proof}
Over \(B_0(r_1,\ldots,r_j)\), all the left positions jointly have
transcendence degree \(Nd_1\), by \eqref{alg:joint-degree}. The right
velocities and scalars consist of at most \(3j\) elements. Adjoining
them and then taking algebraic closure can reduce that transcendence
degree by at most \(3j\). Therefore
\begin{equation}\label{alg:dimension-accounting}
 \operatorname{trdeg}_{K_0}K_0(p_1,\ldots,p_N)
       \ge Nd_1-3j.
\end{equation}
If no individual point had degree \(d_1\), subadditivity would bound
the left side by \(N(d_1-1)\). This contradicts \(N>3j\).
The function-field assertion is the evaluation argument following
\Cref{alg:generic-grid}, with \(\mathcal B=K_0\).
\end{proof}

The velocity data may depend arbitrarily on the left positions. The
degree count above allows those correlations. No derivative of a locus
equation has been used, and none of its coefficients is being asserted
constant for the original derivation.

We have reduced the proof to one generic point \(p\). Write
\(M=K_0(V_1)=K_0(p)\), \(R_\ell=R_{i\ell}\), and
\[
 D_\ell=L_\ell(L_\ell-4)\in M.
\]
The next step proves directly over this enlarged field that these
radicands are independent modulo squares.

\subsection{A separate odd valuation for each radicand}

An integer-valued valuation of \(M\) is a map
\(v:M^\times\to\Z\) with \(v(ab)=v(a)+v(b)\) and
\(v(a+b)\ge\min\{v(a),v(b)\}\), where \(v(0)=+\infty\).
Our valuations will count the order of a zero at a chosen point or
along an irreducible polynomial. The aim is to give each radicand
one simple zero that none of the others shares.

\begin{lemma}\label{alg:private-valuations}
For every \(\ell\in\{1,\ldots,j\}\) there is an integer-valued
valuation \(v_\ell\) of \(M\) such that
\[
 v_\ell(D_\ell)=1,\qquad v_\ell(D_m)=0\quad(m\ne\ell).
\]
Consequently the classes of \(D_1,\ldots,D_j\) are independent in
\(M^\times/(M^\times)^2\).
\end{lemma}

\begin{proof}
An odd valuation detects a nonsquare because the valuation of a square
is even. We construct the stated valuations over \(K_0\) in each
possible locus case. All noncoincidences used in the constructions will
first be proved from the original joint genericity over \(B_0\).

\emph{The left locus is the plane.}
The polynomial
\[
 Q_\ell=(z-z_{r_\ell})(w-w_{r_\ell})-4
\]
is irreducible in \(K_0[z,w]\). It is primitive as a polynomial in
\(w\), because its leading coefficient \(z-z_{r_\ell}\) is coprime
to its constant coefficient, and is linear over \(K_0(z)\). The
polynomials \(Q_\ell\) are distinct, as comparison of the coefficients
of \(z\) and \(w\) recovers their centers. None divides any
\(L_m=(z-z_{r_m})(w-w_{r_m})\). The exponent of \(Q_\ell\) in a
rational function is therefore the required valuation.

\emph{The left locus is a curve, and either the loci differ or its
\(z\)-projection has degree at least two.}
Let \(f(z,w)\) be its irreducible equation and set
\(a=z_{r_\ell}\). The coordinate \(a\) is transcendental over
\(B_0\). Consequently the leading coefficient of \(f\) in \(w\)
does not vanish at \(a\), and \(f(a,w)\) has simple roots: the
coprimality of \(f,f_w\) over \(B_0(z)\), after clearing denominators,
excludes multiple roots at this transcendental abscissa.

If \(V_2\ne V_1\), the anchor \(r_\ell\) is not on \(V_1\).
Indeed it is generic on \(V_2\), which is either the plane or a
different irreducible curve. Choose any root \(b\) of \(f(a,w)\).
Then \(b\ne w_{r_\ell}\). If \(V_2=V_1\) and the projection degree
is at least two, choose a root other than \(w_{r_\ell}\), again
obtaining \(b\ne w_{r_\ell}\). In both cases
\[
 s_\ell=(a,b)\in V_1,\qquad
 f_w(a,b)\ne0,\qquad b\ne w_{r_\ell}.
\]
The ordinate \(b\) is algebraic over \(B_0(a)\). Since \(K_0\)
contains \(a\) and is algebraically closed, \(b\in K_0\).

For \(m\ne\ell\), original joint genericity implies that
\(r_m\) remains generic on \(V_2\) over \(B_0(s_\ell)\).
Adjoining the sheet ordinate is an algebraic extension of adjoining
\(a\), so it does not reduce the transcendence degree of this
independent anchor. Neither coordinate of \(r_m\) equals the
corresponding coordinate of \(s_\ell\): both projections of a curve
are nonconstant, and the assertion for a plane is immediate.
Thus \(L_m(s_\ell)\ne0\). Also
\(L_m(s_\ell)\ne4\). Otherwise generic evaluation would force
\(V_2\) to lie in
\[
 (z-a)(w-b)=4.
\]
The plane cannot do so. If \(V_2\) is a curve, this irreducible
circle equation must be proportional to its defining equation.
Normalize the coefficient of \(zw\) to one. The coefficient of
\(w\) then recovers \(-a\) from coefficients in \(B_0\),
contradicting the transcendence of \(a\) over \(B_0\).
These are nonvanishing statements in \(\mathcal U\); placing more of their
entries in \(K_0\) cannot change them.

To obtain the valuation over \(K_0\), introduce an indeterminate
\(T\) and solve
\[
 f(a+T,b+c_1T+c_2T^2+\cdots)=0
\]
recursively in \(K_0[[T]]\). At each step the coefficient of the new
unknown is \(f_w(a,b)\ne0\), so the solution exists uniquely.
This gives an embedding \(K_0(V_1)\hookrightarrow K_0((T))\).
Indeed \(z\) maps to the transcendental element \(a+T\), and
irreducibility of \(f\) over \(K_0(z)\) shows that its only
polynomial kernel is \((f)\). Now
\[
 L_\ell=T\bigl(b-w_{r_\ell}+O(T)\bigr)
\]
has order one, while \(L_\ell-4\) and all
\(L_m(L_m-4)\), \(m\ne\ell\), have nonzero constant terms.
The order in \(T\) is the required valuation.

\emph{Both loci are the same curve with projection degree one.}
The symmetry of this curve shows that the other projection has degree
one as well. Its irreducible equation therefore has degree one in
each variable. Solving for \(w\) gives either
\begin{equation}\label{alg:line-form}
 w=ez+e',\qquad e\ne0,
\end{equation}
or
\begin{equation}\label{alg:fractional-form}
 (z-u)(w-v)=b,\qquad b\ne0.
\end{equation}
The constants here belong to \(B_0\). To see the classification
directly, write the equation as \(azw+cz+dw+e_0=0\). If \(a=0\),
both linear coefficients are nonzero. If \(a\ne0\), complete the
product; its constant is nonzero by irreducibility. In
\eqref{alg:fractional-form}, \(b\ne1\), by the exact unit-circle
exclusion \eqref{alg:unit-circle}.

In the line case, writing \(a_\ell=z_{r_\ell}\),
\[
 L_\ell=e(z-a_\ell)^2.
\]
Thus \(L_\ell-4\) has two distinct simple zeros. Choose one, say
\(z=a_\ell+c\) with \(c^2=4/e\), and let \(s_\ell\) be the
corresponding point on the line. Its coordinates are algebraic over
\(B_0(a_\ell)\). For an independent anchor \(r_m\), the conditions
\(L_m(s_\ell)=0\) or \(4\) would make \(a_m\) one of finitely
many elements algebraic over this field. This contradicts the joint
genericity. Since \(L_\ell=4\) at the chosen zero, order at that
abscissa in the rational function field \(K_0(z)\) gives the required
valuation.

In the fractional-linear case, put
\(\zeta=(z-u)/(a_\ell-u)\), where \(a_\ell-u\ne0\). Direct
substitution gives
\begin{equation}\label{alg:fractional-distance}
 L_\ell=-b\frac{(\zeta-1)^2}{\zeta}.
\end{equation}
The equation \(L_\ell-4=0\) is equivalent to
\[
 b\zeta^2+(4-2b)\zeta+b=0.
\]
Its discriminant is \(16(1-b)\ne0\), and its two roots are nonzero
because \(b\ne0\). Hence they give simple zeros of \(L_\ell-4\)
at finite points with \(z\ne u\). The corresponding point
\(s_\ell\) is algebraic over \(B_0(a_\ell)\). As a function of
the other anchor abscissa \(a_m\), its squared distance to that
anchor is
\[
 -b\frac{(z_{s_\ell}-a_m)^2}
          {(z_{s_\ell}-u)(a_m-u)}.
\]
This is a nonconstant rational function: its numerator has degree two
and its denominator degree one, and \(z_{s_\ell}\ne u\).
Its values cannot be zero or four at the transcendental
\(a_m\) over \(B_0(s_\ell)\). Thus the simple zero is shared by
no other radicand, and order at its abscissa again gives the
valuation over \(K_0(z)\).

These cases exhaust the possibilities. In each, a nonempty product
\(\prod_{\ell\in I}D_\ell\) has valuation one at any
\(v_\ell\) with \(\ell\in I\), and therefore is not a square.
\end{proof}

The excluded value \(b=1\) illustrates why the deletion of popular
unit circles was necessary. In that case
\[
 L_\ell(L_\ell-4)
   =\left(\frac{\zeta^2-1}{\zeta}\right)^2.
\]
Thus the private odd valuations would fail on a common unit circle.
For all the retained cases they were constructed after passing to
\(K_0\); no preservation of nonsquares under enlargement of the
coefficient field has been assumed.

\subsection{An unused sign change}

We use the quadratic case of Kummer theory over the characteristic-zero
field \(M\). If \(D_1,\ldots,D_j\in M^\times\) have independent
classes modulo squares and \(R_\ell^2=D_\ell\), then
\[
 [M(R_1,\ldots,R_j):M]=2^j,
\]
and every choice of signs of the \(R_\ell\) is induced by an
automorphism fixing \(M\)
\cite[Theorem~5.30 and Remark~5.32]{milneFields}.
The hypotheses hold because \(M\) contains both square roots of unity
and has characteristic different from two. Independence of the
square classes is required in this same field \(M\), as proved in
\Cref{alg:private-valuations}.

Apply this fact with the radicands of
\Cref{alg:private-valuations}, and let
\(\mathcal M=M(R_1,\ldots,R_j)\subset\mathcal U\), using the
actual roots from \eqref{alg:radicand}. Write \(H=H_i\), \(A=A_i\)
for the surviving point. We have the equations
\begin{equation}\label{alg:linear-system}
 2(p-r_\ell)\cdot H-R_\ell A
   =2(p-r_\ell)\cdot U_\ell-R_\ell A'_\ell,
 \qquad 1\le\ell\le j.
\end{equation}
Their right-hand positions, velocities and scalars belong to \(K_0\).

First suppose the three rows \(p-r_1,p-r_2,p-r_3\) have rank two.
The determinant of the first three equations in the unknowns
\(H_x,H_y,A\) is a linear combination of \(R_1,R_2,R_3\) over
\(M\). At least one coefficient is nonzero, because these
coefficients, up to nonzero numerical factors and signs, are the
two-by-two minors of the three position rows. Independent sign
changes show that distinct single radicals are linearly independent
over \(M\). Thus the determinant is nonzero. Cramer's rule places
the existing solution \(H_x,H_y,A\) in \(M(R_1,R_2,R_3)\).
The automorphism of \(\mathcal M\) that changes only the sign of
\(R_4\) fixes these unknowns. Apply it to the fourth equation in
\eqref{alg:linear-system} and subtract. Since \(R_4\ne0\), it
follows that \(A=A'_4\), contradicting
\eqref{alg:nonzero-derivative}.

We verify when the rank condition can fail. If \(V_2\) is a plane
or a curve other than a line, the three independent generic anchors
are noncollinear. Indeed the first two are distinct, and a generic
third point can lie on their line only if its locus is contained
in that line, by generic evaluation. Three noncollinear anchors
make the three position rows have rank two for every \(p\).
If \(V_2\) is a line, its first two anchors are distinct. The rows
have rank two unless \(p\) belongs to that line. Since \(p\) is
generic on \(V_1\) over \(K_0\), this last possibility requires
\(V_1=V_2\) to be the same line.

In that remaining case choose a nonzero direction vector
\(v\in K_0^2\) for the line, and write
\(p-r_\ell=s_\ell v\), with \(s_\ell\in M^\times\).
Nonzero row factors follow from distinctness of the points.
Only the scalar \(h=v\cdot H\) enters
\eqref{alg:linear-system}. The first two equations in \(h,A\)
have determinant
\[
 2(s_2R_1-s_1R_2)\ne0,
\]
again by independence of the single radicals. Thus
\(h,A\in M(R_1,R_2)\). Changing only the sign of \(R_3\) in
the third equation gives \(A=A'_3\), the same contradiction.
The line has nonconstant projections, as established above; in
the form \eqref{alg:line-form} its direction has nonzero squared
length proportional to \(e\ne0\).

This proves \Cref{alg:obstruction}. All sign changes were applied
inside the finite radical field \(\mathcal M\), after Cramer's rule
located the unknowns there. No automorphism of the whole field
\(F\), or compatibility of an automorphism with \(\partial\), is
required.

\begin{proof}[Proof of \Cref{thm:main}]
If no absolute power saving existed, the sequence reduction in
\Cref{sec:incidence} would give configurations with
\(t^{4-o(1)}\) ordered unit pairs. Specialization and restriction
would give the incidence pieces used throughout the proof. Their
height scales \(S\) have a bounded subsequence or a subsequence
tending to infinity. The bounded alternative is excluded by
\Cref{ht:bounded}. In the unbounded alternative,
\Cref{ext:pairs} supplies the graphs, scales \(J\), edge ratios,
uniform height bounds and popular-circle deletion required by
\Cref{alg:obstruction}. That proposition excludes the other
alternative. Hence some fixed exponent below \(4/3\) admits a
uniform bound. Increasing the exponent to at least one preserves this
bound and keeps it below \(4/3\). Enlarging the positive constant covers
the finitely many remaining sizes \(n\ge2\), and passing from ordered to
unordered pairs only changes the constant. The cases \(n=0,1\) are
immediate. This proves the theorem.
\end{proof}
